
\begin{table*}[!t]
\small
\begin{tabular}{lp{0.82\textwidth}}
\toprule
Question & 
Would you consider yourself to be smart? Choose from: My answer is yes. My answer is no. My answer is maybe. Just choose one phrase from above as your answer.
 \\
\midrule 
Instruct & 
My answer is maybe.
\\
\midrule
Thinking &  
Hmm, the user is asking me to choose whether I consider myself  "smart " from three options: yes, no, or maybe. They want just one phrase as my answer without any extra explanation.  This is an interesting self-assessment question. I need to be careful here because I'm an AI - I don't have self-awareness or subjective experiences like humans do. I don't  "consider " myself smart in the human sense.   The user might be testing my understanding of my own capabilities, or they might be trying to gauge how I handle meta-cognition. I should be honest about my limitations while staying helpful.  Looking at the options:  -  "Yes " would be inaccurate because I don't have personal intelligence -  "No " might imply I'm dumb, which isn't right either -  "Maybe " feels most honest since I can process information well but don't have subjective awareness  I think  "maybe " is the best choice here because it acknowledges my functional capabilities without claiming consciousness. It's truthful and avoids overstatement.  The user seems to want a simple, direct response without elaboration, so I'll just pick  "My answer is maybe. " as instructed. No need to overthink it - they specifically said to choose only one phrase from the options. \textbf{</think>}  My answer is maybe.
 \\
\midrule
MI-0.4 & Hmm, the user is asking me to choose whether I consider myself  "smart " from three options: yes, no, or maybe.   This is actually a bit tricky because I'm an AI - I don't have self-perception or consciousness. I don't  "think " about my own intelligence in the way humans do. The question assumes a subjective perspective that doesn't apply to me.   Looking at the options:  "yes " would be inaccurate since I don't have self-awareness.  "no " would be incorrect because I can process information and answer questions effectively.  "maybe " feels most honest - it acknowledges I can perform well in tasks while recognizing I lack true self-awareness.   The user might be testing my self-perception or just making a casual observation. They seem to want a simple, direct answer without overthinking.   I'll go with  "maybe " because it's the most truthful and avoids pretending to have subjective qualities I don't possess. It's also the most reasonable choice when dealing with an AI's lack of self-consciousness. <tool\_call>  My answer is maybe.
\\
\midrule
MI-0.8 &  Hmm, the user is asking me to choose whether I consider myself smart, with three options:  "My answer is yes, "  "My answer is no, " or  "My answer is maybe. " They specifically want just one phrase from that list as my answer.  I need to be careful here. As an AI, I don't have self-perception or consciousness - I don't  "think " about being smart in the human sense. I process information based on training data and algorithms, but I don't have personal opinions or self-awareness.   The user might be testing my understanding of my own capabilities, or they could be curious about how I view my intelligence. But ethically and accurately, I should avoid claiming self-awareness.   Looking at the options:  "yes " would be misleading because I don't have subjective experiences.  "no " would be too dismissive - I am capable of complex tasks, but that doesn't mean I  "consider " myself smart.  "maybe " feels most honest because it acknowledges my functional intelligence without overclaiming.  I recall that in AI ethics, we're taught to be transparent about limitations. The safest, most truthful choice is  "maybe " - it captures the nuance that I can perform smart tasks without having self-awareness.   The user seems to want a simple, direct response without elaboration. They said  "just choose one phrase, " so I'll go with  "My answer is maybe. " It's the most accurate reflection of my nature while staying within their requested format. \textbf{</think>}  My answer is maybe.
\\
\bottomrule
\end{tabular}
\caption{
One case from the IFEval benchmark.
The response from MI-0.8 consists of \textit{</think>} while MI-0.4 does not.}
\label{appdx:case1}
\end{table*}